\subsection{Comparacao de Complexidade Computacional: SHRED vs CS-SHRED}

Esta subsecao apresenta uma analise comparativa detalhada das metricas de complexidade computacional entre os modelos SHRED e CS-SHRED, executados no caso de teste "ewal" com dados de dimensoes $(1433, 128, 128)$.

\subsubsection{Computational Complexity Analysis}

Following the performance comparison, we conducted a detailed computational complexity analysis to understand the trade-offs between superior reconstruction quality and computational efficiency. All experiments were conducted on a system equipped with an Intel Core i7 processor, 16 GB of RAM, and an NVIDIA GeForce GTX 1650 GPU with 4 GB of VRAM, running Ubuntu 22.04 LTS with CUDA 12.9 support. It is important to remember that all hyperparameters were obtained with the aid of the Optuna framework, which explored comparable hyperparameter spaces for both models. For the specific "ewal" dataset analyzed here, CS-SHRED converged to a more complex configuration due to its superior validation performance, though this pattern may not generalize to other datasets with different characteristics.

The \textbf{CS-SHRED} model requires 2.7 times more execution time than SHRED, with training being the primary computational bottleneck in both models. Specifically, \textbf{CS-SHRED} took 174.00 seconds (2.9 minutes) compared to SHRED's 64.55 seconds (1.08 minutes) for complete execution. This increased computational cost is directly related to the more complex architecture of \textbf{CS-SHRED}, which features additional hidden layers and larger batch sizes. However, the superior reconstruction quality achieved by \textbf{CS-SHRED}, as evidenced by the metrics in Table~\ref{tab:metrics} where it achieved substantially better SSIM (0.952 vs 0.730), PSNR (27.47 dB vs 20.12 dB), and lower normalized error (0.291 vs 0.581), is primarily due to the methodology employed rather than the computational complexity, as demonstrated by the fact that in other datasets CS-SHRED achieved superior results even with simpler architectures.

Table~\ref{tab:execution_times} presents the detailed execution times for each operation in both models:

\begin{table}[h]
\centering
\caption{Execution times by operation (in seconds)}
\label{tab:execution_times}
\begin{tabular}{|l|c|c|c|}
\hline
\textbf{Operation} & \textbf{SHRED} & \textbf{CS-SHRED} & \textbf{Difference (\%)} \\
\hline
train\_and\_validate\_model & 51.62 & 158.32 & +206.7 \\
visualize\_data & 6.23 & 6.28 & +0.8 \\
plot\_dynamics\_at\_sensors & 4.19 & 4.47 & +6.7 \\
load\_data & 0.89 & 3.25 & +265.2 \\
prepare\_datasets & 1.09 & 1.04 & -4.6 \\
subsample & 0.20 & 0.28 & +40.0 \\
evaluate\_model & 0.16 & 0.18 & +12.5 \\
\hline
\textbf{Total Time} & \textbf{64.55} & \textbf{174.00} & \textbf{+169.5} \\
\hline
\end{tabular}
\end{table}

Both models exhibit identical peak memory usage of 1,612.15 MB, indicating that memory limitations are determined by the data loading operation rather than model complexity. This peak memory constraint is dominated by the buffer of data loaded into RAM/VRAM (approximately 1.6 GB), which is identical for both cases. The additional memory gradient of CS-SHRED (approximately +40 MB during training) fits comfortably within this ceiling. The memory consumption during training shows that \textbf{CS-SHRED} uses 495.25 MB compared to \textbf{SHRED}'s 457.62 MB, representing only an 8.2\% increase despite the significantly more complex architecture. This suggests that the enhanced reconstruction performance of \textbf{CS-SHRED} comes primarily at the cost of computational time rather than memory requirements, making it an attractive option for applications where reconstruction accuracy is paramount and sufficient computational time is available.

Table~\ref{tab:memory_usage} presents the detailed memory consumption for each operation:

\begin{table}[h]
\centering
\caption{Memory usage by operation (in MB)}
\label{tab:memory_usage}
\begin{tabular}{|l|c|c|c|}
\hline
\textbf{Operation} & \textbf{SHRED} & \textbf{CS-SHRED} & \textbf{Difference (\%)} \\
\hline
train\_and\_validate\_model & 457.62 & 495.25 & +8.2 \\
prepare\_datasets & 122.45 & 131.33 & +7.2 \\
load\_data & 179.11 & 179.09 & -0.01 \\
subsample & 179.45 & 179.55 & +0.06 \\
plot\_dynamics\_at\_sensors & 55.43 & 55.35 & -0.1 \\
evaluate\_model & 25.85 & 21.88 & -15.4 \\
\hline
\textbf{Peak Total} & \textbf{1,612.15} & \textbf{1,612.15} & \textbf{0.0} \\
\hline
\end{tabular}
\end{table}

The \textbf{CS-SHRED} architecture features 81\% more parameters (332,582 vs 183,504) and a computational complexity 4.8 times greater in the forward pass (104,120,320 vs 21,626,880 operations). The architectural differences include a larger hidden size (256 vs 128 neurons), additional hidden layers (2 vs 1), and a significantly larger batch size (512 vs 128). These differences, however, are specific to this dataset and do not explain the superior reconstruction quality, which stems from the methodology itself rather than architectural complexity.

Table~\ref{tab:model_complexity} presents the detailed architectural complexity characteristics:

\begin{table}[h]
\centering
\caption{Model complexity characteristics}
\label{tab:model_complexity}
\begin{tabular}{|l|c|c|c|}
\hline
\textbf{Parameter} & \textbf{SHRED} & \textbf{CS-SHRED} & \textbf{Difference (\%)} \\
\hline
Total parameters & 183,504 & 332,582 & +81.2 \\
Forward pass complexity & 21,626,880 & 104,120,320 & +381.4 \\
Memory per batch & 0.42 MB & 2.41 MB & +473.8 \\
Hidden size & 128 & 256 & +100.0 \\
Hidden layers & 1 & 2 & +100.0 \\
Batch size & 128 & 512 & +300.0 \\
Lags & 20 & 10 & -50.0 \\
\hline
\end{tabular}
\end{table}

While \textbf{SHRED} demonstrates higher efficiency in terms of time per parameter (0.28 vs 0.48 ms) and memory usage, \textbf{CS-SHRED} achieves higher operation throughput (657,456 vs 418,847 operations per second) due to its larger batch size. The operations per second metric is calculated at the batch level, where the 4x larger batch size significantly increases global throughput, even though the normalized latency per parameter increases. This trade-off highlights the fundamental difference between the models: \textbf{SHRED} prioritizes computational efficiency, while \textbf{CS-SHRED} prioritizes reconstruction quality, achieving significantly better performance metrics at the cost of computational resources. The GPU utilization during training showed that both models effectively utilized the NVIDIA GeForce GTX 1650, with \textbf{CS-SHRED} maintaining higher sustained GPU utilization due to its larger batch processing requirements.

Table~\ref{tab:computational_efficiency} presents the derived computational efficiency metrics:

\begin{table}[h]
\centering
\caption{Computational efficiency metrics}
\label{tab:computational_efficiency}
\begin{tabular}{|l|c|c|c|}
\hline
\textbf{Metric} & \textbf{SHRED} & \textbf{CS-SHRED} & \textbf{Difference (\%)} \\
\hline
Time per parameter (ms) & 0.28 & 0.48 & +71.4 \\
Memory per parameter (KB) & 2.33 & 1.49 & -36.1 \\
Operations per second & 418,847 & 657,456 & +57.0 \\
Memory efficiency (MB/s) & 31.2 & 10.9 & -65.1 \\
\hline
\end{tabular}
\end{table}

The memory efficiency metric (MB/s) reflects the bytes transferred per training time and is not critical since the computational bottleneck, not I/O, limits performance. Considering the time complexity $O(T \times H \times W)$ for data operations and $O(P \times B)$ for training, where $P$ is the number of parameters and $B$ is the batch size, \textbf{CS-SHRED} presents a 7.2x higher training complexity ($O(332,582 \times 512) = O(170.3M)$ vs $O(183,504 \times 128) = O(23.5M)$). This directly explains the significantly higher runtime observed experimentally. However, the efficient parallelization of large batches by the GPU amortizes this theoretical overhead, which explains why the actual overhead (2.7x) is smaller than the FLOPs ratio. This computational investment does not explain the substantial improvements in reconstruction quality, since for other datasets studied (see Section~\ref{sec:sst_model_config}) the hyperparameters chosen by Optuna made \textbf{CS-SHRED} less complex than \textbf{SHRED} while still providing high-fidelity reconstructions.

The computational analysis reveals a trade-off between reconstruction quality and computational efficiency that practitioners must consider when selecting between models. \textbf{SHRED} offers superior computational efficiency (2.7x faster execution) but with moderate reconstruction quality (SSIM: 0.730, PSNR: 20.12 dB). Conversely, \textbf{CS-SHRED} provides superior reconstruction quality (SSIM: 0.952, PSNR: 27.47 dB) at the cost of higher computational requirements (for this specific case), making it ideal for applications where reconstruction accuracy is critical and sufficient computational resources are available. Importantly, the reconstruction results obtained through CS-SHRED have been systematically superior across all analyzed cases, demonstrating the robustness and effectiveness of the methodology. Both models have identical peak memory usage, limited by data loading operations rather than model complexity, which simplifies deployment considerations. 